% AA-33 | Audyt I, 22 XI 2026; kontrola matematyczna i redakcyjna.
\chapter{Wypadkowa wielomianów i macierz Sylvestera}\label{aa:ch:resultant}

\section{Motywacja i definicja}
Dla dwóch równań wielomianowych $f(x)=0$ i $g(x)=0$ chcemy
rozstrzygać istnienie wspólnego rozwiązania bez obliczania pierwiastków.
Nad ciałem odpowiedzi udziela największy wspólny dzielnik, lecz
jego obliczenie wymaga dzielenia wielomianów. Wypadkowa jest natomiast
\emph{wyznacznikiem} macierzy zbudowanej z samych współczynników;
można ją więc obliczać również nad pierścieniami innymi niż ciała.

\begin{definicja}\label{aa:def:sylvester}
Niech $A$ będzie pierścieniem przemiennym z jedynką oraz
$m,n\ge1$ liczbami całkowitymi. Ustalmy dwa ciągi
$a_0,\ldots,a_m,b_0,\ldots,b_n\in A$ i wielomiany
$f=\sum_{i=0}^m a_iX^i$, $g=\sum_{j=0}^n b_jX^j$.
Na tym etapie $m,n$ są \emph{stopniami formalnymi}: wolno mieć
$a_m=0$ albo $b_n=0$. Przez $A[X]_{<k}$ rozumiemy wolny
$A$-moduł wielomianów o stopniu mniejszym niż $k$ (łącznie z zerowym),
z bazą $(1,X,\ldots,X^{k-1})$.
\emph{Macierz Sylvestera} $S_{m,n}(f,g)$ jest macierzą homomorfizmu
\[
 \Phi_{f,g}:A[X]_{<n}\oplus A[X]_{<m}\longrightarrow A[X]_{<m+n},
 \qquad (u,v)\longmapsto uf+vg,
\]
w bazach: najpierw współczynniki $u$, następnie $v$, a w przeciwdziedzinie
współczynniki rosnących potęg $X$.
Definiujemy
\[
 \operatorname{Res}(f,g)=\det S_{m,n}(f,g),\qquad
 R_{m,n}(f,g)=(-1)^{mn}\operatorname{Res}(f,g).
\]
W tym rozdziale symbol $R(f,g)$ bez indeksów oznacza $R_{m,n}(f,g)$,
gdy $m=\deg f$, $n=\deg g$ są rzeczywistymi stopniami dodatnimi.
\end{definicja}

\begin{uwaga}\label{aa:uw:res-sign}
Dla $f=a_1X+a_0$ i $g=b_1X+b_0$ otrzymujemy
\[
 S_{1,1}(f,g)=\begin{pmatrix}a_0&b_0\\a_1&b_1\end{pmatrix},\quad
 \operatorname{Res}(f,g)=a_0b_1-a_1b_0,\quad
 R(f,g)=a_1b_0-a_0b_1.
\]
Zatem $R$ odpowiada konwencji $a_1g(-a_0/a_1)$, gdy $a_1$
jest odwracalne. Nazwa \emph{wypadkowa} będzie dalej odnosiła się do $R$.
Formalne dopisywanie zerowych współczynników wiodących może zmienić
wartość wyznacznika i czynnik znakowy. Na przykład dla
$f=2X+1$ i $g=X+3$ nad $\mathbb Q$ mamy $R_{1,1}(f,g)=5$,
natomiast przy formalnym stopniu $2$ wielomianu $g$ otrzymujemy
$R_{1,2}(f,g)=10$, zaś przy formalnym stopniu $2$ wielomianu
$f$ mamy $R_{2,1}(f,g)=-5$. W kryterium wspólnego czynnika stosujemy więc
rzeczywiste stopnie $m,n$; inny wybór stosujemy jawnie przy pochodnej.
Ponadto dla homomorfizmu pierścieni $\rho:A\to B$ zachodzi
$\rho(R_{m,n}(f,g))=R_{m,n}(\rho(f),\rho(g))$ przy
zachowaniu \emph{stopni formalnych}, ponieważ wyznacznik jest
wielomianem o współczynnikach całkowitych w wyrazach macierzy.
\end{uwaga}

\section{Wspólny czynnik i wspólny pierwiastek}
\begin{twierdzenie}[Kryterium wspólnego czynnika]\label{aa:tw:res-common}
Niech $K$ będzie ciałem, a $f,g\in K[X]$ niezerowymi wielomianami
stopni odpowiednio $m,n\ge1$. Przez $\gcd(f,g)$ rozumiemy
moniczny największy wspólny dzielnik. Następujące warunki są równoważne:
\begin{enumerate}
\item $R(f,g)=0$;
\item $\gcd(f,g)$ ma dodatni stopień;
\item $f$ i $g$ mają wspólny pierwiastek w pewnym rozszerzeniu algebraicznym $K$.
\end{enumerate}
\end{twierdzenie}
\begin{proof}
Macierz Sylvestera jest kwadratowa rozmiaru $m+n$. Ponieważ $K$
jest ciałem, jej wyznacznik znika dokładnie wtedy, gdy
$\Phi_{f,g}$ ma niezerowy element w jądrze. To znaczy, że
istnieją $u\in K[X]_{<n}$ i $v\in K[X]_{<m}$, nieoba zerowe,
takie że $uf+vg=0$.

Jeżeli $d=\gcd(f,g)$ ma stopień $r\ge1$, napiszmy $f=df_1$,
$g=dg_1$. Podstawiając $u=g_1$ i $v=-f_1$, otrzymujemy
$uf+vg=0$ oraz $\deg u=n-r<n$ i $\deg v=m-r<m$.
Ponieważ $f,g$ są niezerowe, para $(u,v)$ jest niezerowa,
więc wyznacznik znika.

W przeciwnym kierunku załóżmy, że $\gcd(f,g)=1$ i
$uf=-vg$. Z lematu Euklidesa w pierścieniu $K[X]$ wynika,
że $f\mid v$. Ponieważ $\deg v<m=\deg f$, mamy $v=0$.
Zatem $uf=0$, co w pierścieniu bez dzielników zera oznacza
$u=0$. Jądro jest zerowe, wyznacznik nie znika.

Na koniec w domknięciu algebraicznym $\overline K$ każdy
nierozkładalny czynnik wielomianu ma pierwiastek. Wspólny
czynnik dodatniego stopnia dostarcza wspólnego pierwiastka.
Odwrotnie, jeśli $f(\alpha)=g(\alpha)=0$ dla $\alpha\in\overline K$,
minimalny wielomian $\alpha$ nad $K$ dzieli zarówno $f$, jak i $g$.
\end{proof}

\section{Wzór przez pierwiastki bez założenia rozdzielczości}
\begin{twierdzenie}[Wzór przez pierwiastki]\label{aa:tw:res-roots}
Niech $K$ będzie ciałem, a $f,g\in K[X]$ wielomianami o stopniach
$m,n\ge1$ i niezerowych współczynnikach wiodących $a_m,b_n$.
Jeżeli w ciele rozkładu $L$ wielomianów $f$ i $g$ mamy
\[
 f=a_m\prod_{i=1}^m(X-\alpha_i),\qquad
 g=b_n\prod_{j=1}^n(X-\beta_j),
\]
gdzie pierwiastki powtarzamy zgodnie z ich krotnościami, to
\[
 R(f,g)=a_m^n\prod_{i=1}^m g(\alpha_i)
 =a_m^nb_n^m\prod_{i=1}^m\prod_{j=1}^n(\alpha_i-\beta_j).
\]
Nie zakładamy, że $f$ lub $g$ są rozdzielcze.
\end{twierdzenie}
\begin{proof}
Rozszerzenie skalarów z $K$ do $L$ nie zmienia wyznacznika,
więc możemy pracować nad $L$. Przyjmijmy następującą kolejność
kroków, w których kontrolujemy również znak i współczynnik wiodący.

\emph{Krok 1: odwrócenie bloków.}
W dziedzinie $\Phi_{f,g}$ przestawiamy bloki tak, aby pierwszym
argumentem było $v\in L[X]_{<m}$, a drugim $u\in L[X]_{<n}$.
Przestawienie bloków długości $m$ i $n$ ma wyznacznik
$(-1)^{mn}$, dlatego macierz $\Psi(v,u)=vg+uf$ ma
wyznacznik $R(f,g)$.

\emph{Krok 2: dzielenie przez $f$.}
Każdy $h\in L[X]_{<m+n}$ daje się jednoznacznie zapisać
$h=r+qf$, gdzie $\deg r<m$ i $\deg q<n$.
Ta zmiana współrzędnych jest liniowym izomorfizmem
\[
 T:L[X]_{<m+n}\longrightarrow L[X]_{<m}\oplus L[X]_{<n},
 \qquad h\longmapsto(r,q).
\]
W bazach rosnących potęg macierz odwrotna do $T$ jest blokowo
trójkątna: kolumny $1,\ldots,X^{m-1}$ oraz
$f,Xf,\ldots,X^{n-1}f$ mają wyznacznik $a_m^n$.
Zatem $\det T=a_m^{-n}$.

\emph{Krok 3: działanie na ilorazie.}
Niech $M_g:L[X]/(f)\to L[X]/(f)$ oznacza mnożenie przez klasę
$g$. Dla $\Psi(v,u)=vg+uf$ reszta z dzielenia przez $f$
zależy tylko od $v$ i jest reprezentantem $M_g([v])$.
Współczynnik ilorazu przy $f$ ma postać $u+Q(v)$ dla pewnego
liniowego $Q$. Macierz $T\circ\Psi$ ma więc postać
$\begin{psmallmatrix}M_g&0\\Q&I_n\end{psmallmatrix}$.
Jej wyznacznik jest $\det M_g$. Z drugiej strony jest równy
$a_m^{-n}R(f,g)$, co daje
$R(f,g)=a_m^n\det M_g$.

\emph{Krok 4: uwzględnienie pierwiastków wielokrotnych.}
Niech $P_0=1$ oraz $P_k=\prod_{i=1}^k(X-\alpha_i)$
dla $1\le k<m$. Ponieważ $P_k$ jest moniczny stopnia $k$,
klasy $[P_0],\ldots,[P_{m-1}]$ tworzą bazę $L[X]/(f)$.
W tej bazie mnożenie przez $X$ spełnia
\[
 X[P_{k-1}]=[P_k]+\alpha_k[P_{k-1}]\quad(1\le k<m),
 \qquad X[P_{m-1}]=\alpha_m[P_{m-1}],
\]
bo $P_m=f/a_m$ jest zerem w ilorazie.
Macierz $M_X$ jest zatem trójkątna o przekątnej
$\alpha_1,\ldots,\alpha_m$, nawet jeśli pierwiastki się powtarzają.
Macierz $M_g=g(M_X)$ jest trójkątna o przekątnej
$g(\alpha_1),\ldots,g(\alpha_m)$, więc
$\det M_g=\prod_i g(\alpha_i)$. Daje to pierwszą równość;
druga wynika z podstawienia rozkładu $g$.
\end{proof}

\begin{wniosek}[Symetria wypadkowej]\label{aa:wn:res-sym}
Niech $K$ będzie ciałem, a $f,g\in K[X]$ niezerowymi wielomianami
stopni odpowiednio $m,n\ge1$. Wówczas
$R(g,f)=(-1)^{mn}R(f,g)$.
\end{wniosek}
\begin{proof}
We wzorze z Twierdzenia~\ref{aa:tw:res-roots} oba wyrażenia
zawierają czynnik $a_m^nb_n^m$. Przy zamianie kolejności
pierwiastków każdy z $mn$ czynników $\alpha_i-\beta_j$
zmienia znak. Stąd dokładnie czynnik $(-1)^{mn}$.
\end{proof}

\section{Pochodna formalna i wyróżnik}
\begin{twierdzenie}[Wyróżnik jako wypadkowa]\label{aa:tw:disc-resultant}
Niech $K$ będzie dowolnym ciałem, a $f\in K[X]$ wielomianem
monicznym stopnia $n\ge2$. Wówczas
\[
 \operatorname{disc}(f)=(-1)^{n(n-1)/2}R_{n,n-1}(f,f').
\]
Stopień $n-1$ przy $f'$ jest tutaj formalny: dopuszczamy,
że rzeczywisty stopień pochodnej jest mniejszy lub że $f'=0$.
\end{twierdzenie}
\begin{proof}
Uzasadnimy użycie stopnia formalnego. Kroki 1--3 dowodu
Twierdzenia~\ref{aa:tw:res-roots} wymagają wyłącznie, aby
$f$ miał rzeczywisty stopień $n$ i niezerowy współczynnik wiodący,
a drugi wielomian miał stopień \emph{co najwyżej równy}
zadanemu stopniowi formalnemu $r\ge1$. Istotnie, przestrzeń $u$ ma wtedy
wymiar $r$, dzielenie przez $f$ nadal daje blok jednostkowy
rozmiaru $r$, a ostatni krok dowodu używa jedynie mnożenia przez
klasę drugiego wielomianu. Wobec tego przy $r=n-1$ i monicznym
$f$ otrzymujemy, także gdy $f'=0$, równość
$R_{n,n-1}(f,f')=\prod_{i=1}^n f'(\alpha_i)$.
Jeśli $f=\prod_{i=1}^n(X-\alpha_i)$, to z reguły Leibniza
\[
 f'(\alpha_i)=\prod_{j\ne i}(\alpha_i-\alpha_j),
\]
równość ta pozostaje poprawna również przy powtarzających się
pierwiastkach. W iloczynie po wszystkich $i$ każda para
$i<j$ daje czynnik
$(\alpha_i-\alpha_j)(\alpha_j-\alpha_i)
=-(\alpha_i-\alpha_j)^2$. Par jest $n(n-1)/2$.
To dokładnie definicja wyróżnika, także w dodatniej charakterystyce.
\end{proof}

\begin{przyklad}[Eliminacja zmiennej bez wyznaczania pierwiastków]\label{aa:prz:res-elimination}
Niech $A=\mathbb Z[s,t]$, $f(X)=X^2-s$ i $g(X)=X-t$.
W ciele ułamków $\mathbb Q(s,t)$ mamy $R(g,f)=f(t)=t^2-s$,
a symetria wypadkowej, ponieważ $2\cdot1$ jest parzyste,
daje $R(f,g)=t^2-s$. Możemy też obliczyć wyznacznik
bez przechodzenia do ciała:
\[
 S_{2,1}(f,g)=
 \begin{pmatrix}-s&-t&0\\0&1&-t\\1&0&1\end{pmatrix},
 \qquad\det S_{2,1}(f,g)=t^2-s.
\]
Aby poprawnie odczytać warunek istnienia wspólnego pierwiastka,
wybierzmy dowolne ciało $F$ i dowolne wartości $s_0,t_0\in F$.
Po specjalizacji $s\mapsto s_0$, $t\mapsto t_0$ para równań
$x^2-s_0=0$, $x-t_0=0$ ma wspólne rozwiązanie w rozszerzeniu $F$
dokładnie wtedy, gdy $s_0=t_0^2$: jedynym kandydatem jest $x=t_0$.
Nie trzeba wybierać pierwiastka z $s_0$. Zgodność tej specjalizacji
z obliczonym wyznacznikiem wynika również z uwagi o homomorfizmach
pierścieni po Definicji~\ref{aa:def:sylvester}.
\end{przyklad}

\begin{przyklad}[Pierwiastek wielokrotny]\label{aa:prz:res-repeated}
Niech $K=\mathbb Q$, $f=(X-1)^2$ i $g=X+2$.
Twierdzenie~\ref{aa:tw:res-roots} daje
$R(f,g)=g(1)^2=9$; jest to prawidłowy wynik mimo
nierozdzielczości $f$. Dla pochodnej $f'=2(X-1)$ mamy
$R(f,f')=f'(1)^2=0$ i $\operatorname{disc}(f)=0$.
Pierwsza wypadkowa nie wykrywa wielokrotności pierwiastka
$f$, lecz drugie obliczenie już tak.
\end{przyklad}

\begin{przyklad}[Kontrprzykład nad pierścieniem z dzielnikami zera]\label{aa:prz:res-zero-divisors}
Niech $A=\mathbb Z/4\mathbb Z$ i $f=g=2X+1$.
Wyznacznik macierzy Sylvestera jest równy zero, ponieważ
jej obie kolumny są identyczne. Jednak żaden $x\in A$
nie spełnia $f(x)=0$: wartość $2x+1$ jest zawsze klasą
nieparzystą. Pierścień $A$ nie jest ciałem, a jego niezerowy
współczynnik wiodący $2$ nie jest odwracalny.
Co więcej, $f$ jest elementem odwracalnym $A[X]$, ponieważ
$(1+2X)(1-2X)=1-4X^2=1$. Nie ma więc pierwiastka nawet
w żadnej niezerowej przemiennej $A$-algebrze: obraz elementu
odwracalnego pozostaje odwracalny. Nie wolno bez dodatkowych
założeń przenosić równoważności z
Twierdzenia~\ref{aa:tw:res-common} na dowolne pierścienie.
\end{przyklad}
